% =============================================================================
% INTERCALATIO: DE ARBORE ET VOLATU
% Status: DRAFTED
% Placement: Mid-book, after system has destabilized but before transmission
% Pure observation — no characters, no dialogue, no typographic distortion.
% A bird murmuration as structural meditation. Returns to pre-narrative mode
% but now the reader reads it differently. The stability is retrospective.
% =============================================================================

\chapter*{Intercalatio: De Arbore et Volatu}
\addcontentsline{toc}{chapter}{De Arbore et Volatu}

\setstretch{1.05}

\begin{adjustwidth}{2.4cm}{2.4cm}

The tree stands slightly apart from the others, not isolated, but with enough
space around it that its full shape can be taken in at once. Its trunk rises with
a gradual, almost hesitant widening at the base, as if it had once adjusted itself
to the ground and then decided to remain. The bark is not uniform. In some places
it lies flat and close, smooth enough to reflect a dull sheen of light, while in
others it lifts in thin plates, curling at the edges, exposing a darker layer beneath.
There are shallow fissures running vertically, not deep enough to divide, but enough
to suggest a slow and continuous separation that never completes.

The branches begin higher than expected. For several feet the trunk is uninterrupted,
a column that holds itself without deviation, and then, almost all at once, the
structure opens. Limbs extend outward in uneven directions, not symmetrically, but
with a balance that only becomes apparent when seen as a whole.

Leaves gather densely toward the outer portions of the canopy. Up close, they are
small and irregular, each one slightly different in shape, but from a distance they
merge into a continuous surface that filters the light into a shifting pattern below.
The wind moves through them without urgency. It does not pass evenly. Certain
clusters respond first, trembling in place, while others remain still for a moment
longer, and then follow, as though the motion were being transmitted through an
unseen delay.

The ground beneath the tree is worn, though not barren.

\vspace{2em}

For a while, nothing happens.

\vspace{2em}

Then, without warning or preparation, the first bird arrives.

It does not descend gradually. It appears in motion already committed, wings
adjusting only at the last moment, the body angling to meet a particular branch.
The landing is precise but not delicate. The branch dips slightly under the added
weight, then returns, oscillating once before settling. The bird holds still, head
turning once, quickly, as if verifying its position.

A second bird follows, then a third, not in a coordinated pattern, but not entirely
independent either. They arrive from different directions, but within a narrow span
of time that suggests a shared decision made elsewhere.

Soon there are many.

They do not fill the tree evenly. Certain branches collect them, forming dense lines
of bodies facing roughly the same direction, while other branches remain empty, not
avoided, but not chosen. Their motion does not cease after landing. It changes form.

The sound begins after the settling.

At first it is intermittent, individual calls that do not align. Then it thickens, not
into a unified chorus, but into a field of overlapping signals. No single call dominates,
yet none are entirely lost. The result is not noise, exactly, but something structured
at a level that resists immediate parsing.

One bird leaves.

Its departure is abrupt, the inverse of its arrival. It pushes off, wings opening into
the space it had just occupied, and is gone before the surrounding birds have fully
responded. The branch lifts slightly in its absence. For a moment, there is a small
gap where it had been, a visible interruption in the line.

Another bird moves into that space.

Not immediately, and not by direct intention that can be observed, but the gap does
not remain.

\vspace{2em}

The birds remain for a time that is neither brief nor extended enough to establish a
rhythm. Then, as without warning as they arrived, they begin to leave.

Not all at once.

One, then another, then several in quick succession. The branches lighten
incrementally. The sound thins. The structure that had formed dissolves without
reversing itself. There is no final departure, no last bird that marks the end.

At some point, the tree is simply empty again.

It is not the same as before.

Nothing visible has changed. The bark is as it was. The branches hold their
positions. The leaves continue to move with the wind.

And yet the space retains a trace of the arrangement that occupied it, not as a
visible residue, but as a condition that could recur.

The tree stands.

It does not anticipate.

It does not remember.

But the possibility remains, suspended in the same way the leaves hold light,
waiting for nothing in particular, and for that reason, available to everything that
might arrive.

\vspace{2em}

\begin{center}
Drip.
\end{center}

\end{adjustwidth}
